
% =============================================================================
\section{Remediation and Incident Response}
\label{sec:remediation}

\textbf{Primary fix: upgrade to ZCS 10.1.20} (2026-07-20). The one-function
rewrite (Section~\ref{sec:vuln}) moves the attacker-controlled value across
the \texttt{execve} boundary; our A/B lab result
(Section~\ref{sec:fix}) confirms the sink goes inert while the positive
control still fires. Because CISA KEV lists the CVE with a 2026-08-24
federal deadline and CERT Polska documents active exploitation, we treat
``upgrade + deploy detections'' as the minimum bar for any internet-adjacent
ZCS deployment.

\textbf{Interim mitigations} (if the upgrade is delayed), in preference
order: (1) \texttt{zmlocalconfig -e snmp\_notify=0} + SNMP service restart
(disarms the trigger), (2) stop \texttt{zimbra-snmp} entirely if SNMP
monitoring is unused, (3) remove the mail-driven \texttt{watchfor} block from
\texttt{swatchrc.in}/\texttt{swatchrc}, (4) tighten the SMTP attack path
(\texttt{mynetworks} to localhost, mandatory authenticated submission).
None of these replaces the upgrade; they shrink the window. Full command
detail: \texttt{remediation.md}.

\textbf{Incident response.} \texttt{playbook-ir.md} provides a four-phase
playbook validated against the lab artifacts: (Phase 0) read-only triage —
grep \texttt{zmswatch.out} for \texttt{SNMP notification:} lines with shell
syntax, grep the mail logs for \texttt{to=<\$(,} capture the live process
tree (the \texttt{ps} shape of Listing in Section~\ref{sec:detection}),
md5-compare the login banners against baseline, hunt zimbra-owned recent
files and cron/SSH persistence; (Phase 1) containment — kill the captured
chain, disable \texttt{snmp\_notify}, block attacker egress, preserve a
hashed evidence tarball before eradication; (Phase 2) eradication — restore
stock web assets from a clean build, purge the poisoned queue entries,
rebuild if residue is unexplainable; (Phase 3) recovery — patch first,
re-verify with the canary detector, keep the detection package deployed,
30-day enhanced monitoring.

\textbf{The defender's recurring question} — ``how do I know I'm not
already hit?'' — has a concrete answer now: run the passive
\texttt{--audit} today, review mail logs for the \texttt{to=<\$(} pattern
over the last 90 days (our YARA rule does this at 166\,MB scale), compare
banner hashes, and check for the process shape of
Section~\ref{sec:detection} in EDR history.

% =============================================================================
\section{Ethical Considerations, Reproducibility, and the Agent Pipeline}
\label{sec:ethics}

\textbf{Responsible disclosure.} CVE-2026-73570 is a publicly disclosed
vulnerability (vendor advisory 2026-06-26; CERT Polska 145/2026; CISA KEV
2026-08-21). We make no discovery claim; this work is independent lab
verification, empirical trigger identification, and defender tooling. All
exploitation was performed exclusively inside the isolated, zero-egress
Docker lab of Section~\ref{sec:lab} against a fictional internal domain
(\texttt{mail.zmdemo2.lab}); no third-party system was touched. Payloads
followed minimum-impact discipline: canary-first (a file timestamp is the
only effect of the proof payload), one defaced asset per service with
backups and documented restoration, and full lab restore verified by hash
after every experiment.

\textbf{AI agents in the pipeline.} This research and its lab environment
were built and operated with AI agents in a controlled environment.
Beyond that note, we are specific about what the agent did, because the
distinction matters: the agent executed the full operational pipeline —
recon, trigger-matrix testing, dry-run/kill-shot sequencing, process and
log forensics, defacement and its verification, lab restoration — under a
hard think-then-act barrier, with every decision logged to the beat log
that drives Figure~\ref{fig:agent}. A human supervised scope, approved the
kill-shot, and verified each on-camera claim frame-by-frame.

This gives an empirical data point for the LLM-pentesting literature.
CHECKMATE~\cite{wang2025checkmate} identifies incoherent long-horizon
planning as the dominant LLM-agent failure mode and proposes an external
classical planner. Our 7-stage beat log (RECON$\rightarrow$\dots$\rightarrow$
FLAG, 51 timed beats in the cinematic run, zero stage regressions, zero
re-plans mid-exploit) is a real-vulnerability instance where the coherence
problem did not surface — consistent with the argument that
\emph{structural} barriers (think/act separation, an explicit stage
artifact, verifiable done-criteria per stage) substitute for, or complement,
an external planner. We report this as an n=1 case study, not evidence of
general agent capability; the beat log is released so others can check it.

\textbf{Reproducibility.} The companion folder
\texttt{Research/CVE-2026-73570/} is self-contained: \texttt{lab/}
rebuilds the target in one script (image build $\approx$15--25\,min with
provisioning egress; runtime stays 0-egress); \texttt{exploit/} and
\texttt{detection/} are stdlib-only Python; \texttt{rules/} imports
directly into Sigma-compatible SIEMs, YARA 4.x, and Suricata;
\texttt{ioc/}, \texttt{remediation.md}, and \texttt{playbook-ir.md} carry
the defensive artifacts; \texttt{evidence/} (with SHA-256 manifests) holds
every figure's source frame and the live-fire transcripts cited in this
paper; \texttt{videos/} holds the two approved recordings. A defender can
go from clone to ``my detections catch this attack'' in under an hour.

% =============================================================================
\section{Conclusion}
\label{sec:conclusion}

CVE-2026-73570 is a textbook trust-chain break: an MTA that logs attacker
text, a monitor that captures it, and a sink that executes it — three
subsystems each doing something reasonable, together doing something
lethal. The public record named the sink; this paper supplies the middle
and the back end: the empirically identified recipient-field trigger path
with its wire constraints, negative-control proof of execution, live
process-tree evidence, a fix verified by A/B lab comparison, and a
defender package (5 Sigma / 3 YARA / 3 Suricata + EDR mappings + canary
self-test + IR playbook) in which every detection was shown firing against
the real attack. The same pipeline was operated end-to-end by autonomous AI
agents under a logged think-then-act discipline — a case study in
long-horizon attack coherence that we release in full, evidence included,
because the defenders who need it most are the ones who can least afford
to trust vendor summaries alone.

% =============================================================================
\begin{thebibliography}{99}

\bibitem{wang2025checkmate}
L.~Wang, X.~Shi, Z.~Li, Y.~Jiang, S.~Tan, Y.~Jiang, J.~Cheng, W.~Chen,
X.~Shen, Z.~Li, and Y.~Chen, ``Automated Penetration Testing with LLM Agents
and Classical Planning,'' \emph{arXiv:2512.11143 [cs.CR]}, 2025.

\bibitem{nvd73570}
NIST NVD, ``CVE-2026-73570: Command injection in the SNMP monitoring
component of Zimbra Collaboration Suite before 10.1.20,''
\url{https://nvd.nist.gov/vuln/detail/CVE-2026-73570}.

\bibitem{zimbraadvisory}
Zimbra, ``Zimbra Security Advisories — command injection in SNMP
monitoring component (advisory 2026-06-26; fixed in 10.1.20, 2026-07-20),''
\url{https://wiki.zimbra.com/wiki/Zimbra_Security_Advisories}.

\bibitem{certpolska}
CERT Polska, ``Advisory 145/2026: actively exploited vulnerability in
Zimbra Collaboration Suite,'' 2026-08-17,
\url{https://moje.cert.pl/komunikaty/2026/145/}.

\bibitem{ciskev}
CISA, ``Known Exploited Vulnerabilities Catalog — CVE-2026-73570
(added 2026-08-21, FCEB due 2026-08-24),''
\url{https://www.cisa.gov/known-exploited-vulnerabilities-catalog}.

\bibitem{zmcommit}
Zimbra, ``zm-build commit c4837c66 —
doSNMP: backtick command to multi-arg system(),''
\href{https://github.com/Zimbra/zm-build/commit/c4837c66360979d757890271f6c976727e71bcad}{\texttt{https://github.com/Zimbra/zm-build/commit/\hspace{0pt}c4837c663609\hspace{0pt}79d757890271\hspace{0pt}f6c976727e71bcad}}.

\bibitem{zimbraadminguide}
Zimbra, ``Admin Guide — SNMP Monitoring,''
\url{https://github.com/Zimbra/adminguide/blob/develop/monitoring.adoc}.

\bibitem{sigma}
SigmaHQ, ``Sigma — generic signature framework for SIEM event logs,''
\url{https://github.com/SigmaHQ/sigma}.

\bibitem{yara}
VirusTotal, ``YARA — the powerful rule-based matching tool,''
\url{https://github.com/VirusTotal/yara}.

\bibitem{suricata}
OISF, ``Suricata — network IDS/IPS/NSM engine,''
\url{https://suricata.io/}.

\end{thebibliography}

\end{document}
